\begin{table}[H]
\centering
\caption{Matriz de confusão do Protocolo B, somada nos 24 \emph{folds} (rodada \texttt{exp015}). Linhas: classe verdadeira do segmento; colunas: previsão do detector binário. Cada segmento de falha é testado seis vezes (uma por bloco normal) e cada segmento normal quatro vezes (uma por falha deixada de fora); as porcentagens são por linha. Total: VP = 1.062, FN = 354, VN = 236, FP = 0.}
\label{tab:confusaoB}
\begin{tabularx}{\textwidth}{>{\raggedright\arraybackslash}X C{2.4cm} C{3.4cm} C{3.4cm}}
\toprule
 & & \multicolumn{2}{c}{\textbf{Previsto}} \\
\cmidrule(lr){3-4}
\textbf{Classe verdadeira} & \textbf{Testes} & \textbf{normal} & \textbf{falha} \\
\midrule
\texttt{normal} & 236 & 236 (100,0\%) & 0 (0,0\%) \\
\midrule
\texttt{bpfi\_0.3mm} & 354 & 0 (0,0\%) & 354 (100,0\%) \\
\texttt{bpfi\_1.0mm} & 354 & 0 (0,0\%) & 354 (100,0\%) \\
\texttt{bpfo\_0.3mm} & 354 & 354 (100,0\%) & 0 (0,0\%) \\
\texttt{bpfo\_1.0mm} & 354 & 0 (0,0\%) & 354 (100,0\%) \\
\bottomrule
\end{tabularx}
\end{table}
